% titlepages/master-title.tex -- титульна сторінка магістерської
% дисертації (Додаток 6 до "Рекомендацій до структури та змісту
% кваліфікаційних робіт здобувачів ступеня бакалавра та магістра",
% КПІ ім. Ігоря Сікорського, 2022).
%
% Відмінності від бакалаврської (bachelor-title.tex): додається гриф
% "На правах рукопису" з УДК зліва вгорі, керівника називають
% "Науковий керівник", а тип освітньої програми (ОПП чи ОНП) береться
% з config.tex (\programmeType), бо для магістра це не завжди ОПП.
\thispagestyle{empty}

\begin{center}
  Національний технічний університет України\\
  «Київський політехнічний інститут імені Ігоря Сікорського»
\end{center}

\vspace{1em}
\noindent\facultyName\\
\departmentName

\vspace{2em}
\noindent
\begin{minipage}[t]{7.5cm}
  «На правах рукопису»\\
  УДК~\udcCode
\end{minipage}%
\hfill
\begin{minipage}[t]{7.5cm}
  \raggedleft
  ДО ЗАХИСТУ ДОПУЩЕНО\\
  Завідувач кафедри\\
  \hrulefill\ \departmentHeadName\\
  «\hrulefill»\hrulefill~\thesisYear~р.
\end{minipage}

\vspace{2em}
\begin{center}
  Магістерська дисертація\\
  на здобуття ступеня магістра\\
  за \programmeType~програмою «\programmeName»\\
  зі спеціальності \specialtyCode~«\specialtyName»\\
  на тему: «\thesisTitle»
\end{center}

\vspace{2em}
\noindent Виконав(-ла):\\
студент(-ка) \authorCourse~курсу, групи \authorGroup\\
\authorFullName \hfill \hrulefill

\medskip
\noindent Науковий керівник:\\
\supervisorRegalia,\\
\supervisorFullName \hfill \hrulefill

\ifdefempty{\consultantFullName}{}{%
  \medskip
  \noindent Консультант з \consultantChapter:\\
  \consultantRegalia,\\
  \consultantFullName \hfill \hrulefill
}

\medskip
\noindent Рецензент:\\
\reviewerRegalia,\\
\reviewerFullName \hfill \hrulefill

\vspace{2em}
\noindent Засвідчую, що у цій магістерській дисертації немає
запозичень з праць інших авторів без відповідних посилань.\\
Студент (-ка) \hrulefill

\vfill
\begin{center}
  Київ~--~\thesisYear
\end{center}

\clearpage
